\section{Interception}\label{sec:intercept}
Interception starts when an object has been detected as a drone. In the
baseline, this is when the detector's confidence exceeds a threshold $\tau$
in $k$ consecutive images. In the target system, the interceptor starts
flying at the first detection by the drone's own model, and the base station
confirms the target or calls the interceptor back while it flies
(Section~\ref{sec:base-station}). The centre of the bounding box is converted
into a direction to the target, using the camera's intrinsic parameters and
the gimbal's orientation, as in \cite{zheng2025mcm}. Two PID controllers turn
the gimbal so that the target stays in the centre of the image, and a Kalman
filter smooths the direction, estimates how fast it turns, and predicts it
when the target is briefly lost. The gimbal loop must react fast enough to
keep the target in view: by Equation~\eqref{eq:angrate}, a target crossing at
17\,m/s leaves a 20$^\circ$ view within about 0.2\,s at 20\,m if the gimbal
lags. This is one reason the target system considers a wide view on the
gimbal for close range (Section~\ref{sec:camera-choice}).

\subsection{Guidance Laws}

Three guidance laws are compared. Pure pursuit and proportional navigation
need only the direction to the target and its rate of change, which one
camera provides; the MPC also needs the target's position and velocity, which
one camera measures poorly.

\textbf{Baseline: pure pursuit.} The interceptor flies at full speed towards
the target's current direction, while the gimbal keeps the target in the
centre of the image. Its weakness is that it always aims at where the target
\emph{is}, not where it is going, so against a crossing or turning target it
ends up behind the target and has to chase it. It can therefore only catch a
target that is slower than itself: against a hostile drone faster than $v_i$
(Table~\ref{tab:assumptions}), it falls behind and never hits.

\textbf{Target system: MPC.} The interceptor uses the MPC of Zheng et al.\
\cite{zheng2025mcm}. It predicts the target's trajectory over $K$ steps from
its estimated position and velocity, assuming constant velocity, and chooses
the accelerations $\mathbf{u}(k)$ that minimise
\begin{equation}
  J = \sum_{k=1}^{K} \|\mathbf{p}_{\mathrm{exp}}(k) - \mathbf{p}(k)\|^2_{Q}
      + \|\mathbf{u}(k)\|^2_{R},
\end{equation}
where $\mathbf{p}_{\mathrm{exp}}$ is the target's predicted position and
$\mathbf{p}$ the interceptor's. The problem has an analytical solution; the
first acceleration is sent to PX4 and the problem is solved again at the next
step. Unlike \cite{zheng2025mcm}, who keep an offset around the target to fire
a net, the offset here is zero. Since the hostile drone does not always fly
in a straight line, the constant-velocity prediction is wrong when it turns,
but it is corrected at every step as new estimates arrive; Zheng et al.\
used the same model against a target flying in a circle. The MPC needs the
target's position and velocity, which one camera cannot give. It is therefore
only possible with triangulation from two or more defenders, for which the
base station sends a second defender (Section~\ref{sec:base-station}), and is part of
the target system.

\subsection{Evaluation}
Metrics are computed by the scenario runner from the true positions in
Gazebo, never from the interceptor's estimates. The target always flies
towards the base, and a run ends with a \emph{hit} (interceptor within
1.5\,m of the target), the \emph{target reaching the base}, the \emph{target
being lost}, or a \emph{timeout}. Head-on, the two drones close at up to
$v_i + v_h \approx 29$\,m/s, about 2.9\,m between two samples at 10\,Hz, which
is more than the hit radius. The closest distance is therefore computed along
the straight line between consecutive samples, not only at the samples.

\textbf{Alternative success: lock-on.} How the hostile drone is finally
stopped is outside the scope of this project, and flying into it is only one
option. An interception may therefore also be counted as successful when the
interceptor holds a \emph{lock-on} on the hostile drone for $t_{\mathrm{lock}}$
seconds. The interceptor has a lock-on when the target is inside its gimbal
camera's view, the detector confirms it as a drone in at least $k$ of every
$n$ consecutive images, and the interceptor is within a distance
$d_{\mathrm{lock}}$ of it. The lock-on must last $t_{\mathrm{lock}}$ seconds
without a break, and be held before the target comes within $x$\,m of its
goal (Section~\ref{sec:evaluation}). A sustained lock-on shows that the
interceptor could follow and engage the target with whatever means is chosen
later. The assumed values of $t_{\mathrm{lock}}$ and $d_{\mathrm{lock}}$ are
listed in Table~\ref{tab:target-assumptions}; both success criteria are
reported for every run.

For each run the following are recorded:

\begin{itemize}
    \item \textbf{Hit or not:} Whether the interceptor successfully intercepted the target.
    \item \textbf{Time to hit:} The time elapsed from first detection to interception.
    \item \textbf{Lock-on:} Whether a lock-on was held for $t_{\mathrm{lock}}$ seconds, the time from first detection until it started, and the longest unbroken lock-on in the run.
    \item \textbf{Distance to base at hit:} How far the target was from the protected asset when neutralized.
    \item \textbf{Target's closest distance to base:} A measure of how close the threat got before being stopped.
    \item \textbf{Closest distance between interceptor and target:} Useful for evaluating near-misses in unsuccessful runs.
\end{itemize}

All methods are run on the same seeded scenarios: a target approaching
head-on and crossing at 30, 60 and 90$^\circ$, at increasing speeds from
4\,m/s up to and beyond the interceptor's top speed $v_i$, including the
hostile speed $v_h$ of Table~\ref{tab:assumptions}, each 30--50 times. Each
method is run both with YOLO detections and with perfect detections from the
true position, which separates errors in guidance from errors in detection.
The measured time to hit is also compared with the margins predicted in
Table~\ref{tab:feasibility}, which tests the analysis in
Section~\ref{sec:derived}. When the target is lost, another drone that still
sees it could take over; this is not evaluated here.

Since the sensors have not yet been tested, realistic distances and speeds
are not known. Zheng et al.\ \cite{zheng2025mcm} captured a target moving at
4~m/s in 64.7\% of their outdoor attempts. The project aims for a similar
share of hits. Pure pursuit is expected to hit only while the target is
slower than $v_i$; proportional navigation and the MPC are expected to hit
more often, earlier and further from the base.